% Generated draft: requires mathematical and source review.
\documentclass[11pt]{article}
\usepackage[margin=25mm]{geometry}
\usepackage{fontspec}
\setmainfont{Latin Modern Roman}
\usepackage{amsmath,amssymb,mathtools,unicode-math}
\setmathfont{Latin Modern Math}
\usepackage{enumitem,needspace,xurl,hyperref}
\setlength{\parindent}{0pt}
\setlength{\parskip}{5pt}
\setlength{\emergencystretch}{3em}
\newcommand{\N}{\mathbb{N}}
\newcommand{\Z}{\mathbb{Z}}
\newcommand{\Q}{\mathbb{Q}}
\newcommand{\R}{\mathbb{R}}
\newcommand{\C}{\mathbb{C}}
\newcommand{\F}{\mathbb{F}}
\newcommand{\A}{\mathbb{A}}
\newcommand{\PP}{\mathbb{P}}
\newcommand{\E}{\mathbb{E}}
\newcommand{\eps}{\varepsilon}
\newcommand{\dd}{\,\mathrm{d}}
\begin{document}
\section*{UnsolvedMath ID: 30005304}
\textbf{Source problem code:} catalogue entry
\subsection{Definitions and mathematical statement}
The supplied passage labels the following assertion Conjecture 3:
\[
\operatorname{supp}(P)\text{ is a lattice},\qquad P\in\mathcal M(G)
\quad\Longrightarrow\quad P\in\mathcal M_F(G).
\]
Here \(\mathcal M(G)\) denotes the source's class of Markov distributions for \(G\), and \(\mathcal M_F(G)\) denotes its class of distributions that factorize over \(G\). The supplied extract does not define these classes, the permitted graphs and state spaces, or the order used for the lattice condition. Consequently, the displayed assertion preserves the source formulation but does not determine a unique fully specified conjecture.

For clarity, one possible precise specialization is the following, conditional on conventions that require confirmation against the source. Let \(G=(V,E)\) be a finite simple undirected graph, let each \(\mathcal X_v\) be a finite totally ordered set, and equip \(\mathcal X=\prod_{v\in V}\mathcal X_v\) with coordinatewise minimum \(x\wedge y\) and maximum \(x\vee y\). For a probability mass function \(p\) on \(\mathcal X\), define
\[
S_p=\{x\in\mathcal X:p(x)>0\}.
\]
Interpret the lattice condition as closure in the ambient product:
\[
x,y\in S_p\quad\Longrightarrow\quad x\wedge y\in S_p\ \text{ and }\ x\vee y\in S_p.
\]
Interpret \(P\in\mathcal M(G)\) as the global Markov property: for every three pairwise disjoint subsets \(A,B,C\subseteq V\), if every path from \(A\) to \(B\) meets \(C\), then \(X_A\) and \(X_B\) are conditionally independent given \(X_C\) under \(P\). Let \(\mathcal C(G)\) be the set of maximal cliques of \(G\). Interpret \(P\in\mathcal M_F(G)\) as the existence of nonnegative functions \(\psi_D:\prod_{v\in D}\mathcal X_v\to[0,\infty)\), for \(D\in\mathcal C(G)\), such that
\[
p(x)=\frac{1}{Z}\prod_{D\in\mathcal C(G)}\psi_D(x_D)
\quad\text{for every }x\in\mathcal X,\qquad
Z=\sum_{x\in\mathcal X}\prod_{D\in\mathcal C(G)}\psi_D(x_D)>0.
\]
Under these conditional conventions, the question is whether every globally Markov \(P\) whose support is closed under coordinatewise minimum and maximum admits this factorization. These finite-state, ordering, Markov, and factorization conventions are an explicit possible interpretation, not hypotheses established by the supplied record.

The source describes Conjecture 3 as a strengthening motivated by total positivity. Total positivity is not an additional hypothesis in the displayed conjecture. The extract refers to an unavailable equation (1) for the observation that totally positive distributions have lattice support; that observation and its scope cannot be checked from the supplied passage.
\subsection{Short English statement}
The question asks whether a Markov distribution must factorize over its graph whenever its support is a lattice. The supplied record labels this conjecture open, but leaves the state spaces, order, Markov property, and factorization conventions unspecified. The finite discrete formulation above is a conditional interpretation requiring human mathematical review.
\subsection{Sources}
\begin{itemize}
\item Local UnsolvedMath registry, supplied record with permanent ID \(30005304\), exact source problem code \texttt{catalogue entry}, title ``Lattice Supports and Graphical Model Factorization,'' in the collection ``Oberwolfach Reports Open Problems.'' The record supplies no named proposer and records proposed year 2022.
\item Supplied report citation: \emph{Algebraic Structures in Statistical Methodology} (2023), Oberwolfach Reports, pp.\ 3121--3170, DOI \texttt{10.4171/owr/2022/55}. Supplied source URL: \url{https://doi.org/10.4171/owr/2022/55}. The extract identifies the factorization assertion as Conjecture 3. This report citation is supplied provenance and has not been independently consulted.
\item Supplied literature source, labeled ``Lauritzen OWR report'': \url{https://doi.org/10.4171/owr/2022/55}. The exact supplied literature-check date is \texttt{2026-08-22}. The supplied assessment labels the problem open and says that the strengthened factorization claim remains unresolved in the stated generality, while the extracted passage subsequently turns to a different optimization discussion. This is supplied dated curation, not an independently verified literature assessment.
\item The extract mentions references [1] concerning conditions for factorization and [4] concerning binary totally positive distributions, and invokes equation (1). Their bibliographic entries and the equation are absent. Reference [5] occurs in the subsequent Gaussian optimization discussion; its bibliographic entry is also absent.
\item The registry labels the statement \texttt{corrected\_verified} and explains that only formula delimiters, notation, grammar, and typography were repaired, with mathematical content and scope preserved. This supplied editorial label does not establish mathematical correctness or resolve the missing definitions.
\end{itemize}
\end{document}
